% !TeX program = XeLaTeX
% !TeX root = ../pujavidhanam.tex
% --meta--
% source: श्रीगङ्गाकल्पोक्तपूजाविधिः, ed. Mittur Srinivasa Bhatta (Idkidu, 2015), pp. 41-44.
% --end-meta--
\chapt{काश्यादि-तीर्थक्षेत्र-यात्रा-प्रस्थान-विधिः}

आचम्य पवित्रपाणिः गुर्वाद्यभिवन्दनं जलशुद्ध्यादिकं च कृत्वा अक्षतान् गृहीत्वा

\dnsub{सङ्कल्पः}

`सुमुखश्चैकदन्तश्च\ldots' शुभतिथौ (नक्षत्र-राशि-नामोच्चारणं कृत्वा)
अस्मत्पितृमातृवंशजातानां दिवि भुवि तिर्यङ्नरकादिषु वासिनाम् अन्येषां च
अस्मत्पूर्वजातानाम् उत्तमलोकावाप्त्यर्थं समस्तपूर्वजोद्धारार्थं नान्दीमुखाश्रुमुखपितृप्रीत्यर्थम्
ऊष्मपो यमोऽङ्गिरस्वानग्निष्वात्तोऽग्निः कव्यवाहनादयो ये पितरस्तेषामविच्छिन्नप्रसादसिद्ध्यर्थम्
असृग्भागि-लेपभागि-पिण्डभागि-विप्रोच्छिष्टभागि-प्रभृतीनाम् अक्षय्यतृप्त्यर्थं
नानापुण्यनदीस्नानपूजाफलसिद्ध्यर्थं मम जन्मप्रभृत्यद्य यावत्
ज्ञाताज्ञात-कामाकाम-सकृदसकृत्कृत-कायिक-वाचिक-मानसिक-सांसर्गिक-स्पृष्टास्पृष्ट-पीतापीत-भुक्ताभुक्त-सकलपातकोपपातक-लघुपातक-सङ्करीकरण-मलिनीकरणापात्रीकरण-जातिभ्रंशकरादीनां
स्वतः परप्रेरणया वा सम्भावितानां सकलविधपापानां परिहारार्थं
यमुना-सरस्वती-चन्द्रभागाख्य-महापुण्यनदीस्नानकल्पोक्तफलसिद्ध्यर्थं
वेत्रवती-मलापहा-फल्गुनी-कृष्णाख्य-पवित्रनदीक्षेत्रयात्राफलसिद्ध्यर्थं
शुतुद्री-गोदावरी-मरुद्वृधाख्य-पवित्रनदीतीर्थमज्जन-शतफलसिद्ध्यर्थं
भीमरथी-सीता-चर्मण्वती-नर्मदाख्य-पुण्यनदीमाहात्म्यश्रवणफलावाप्त्यर्थं
पयोष्णी-विपाशा-ताम्रपर्णीति तत्तन्नदीकल्पोक्त-तत्तत्फलसिद्ध्यर्थं
कावेरी-करतोया-नेत्रावत्यादि-नदीविशेष-तीर्थविशेष-स्नानतर्पणफलप्राप्त्यर्थं
गङ्गाविश्वेश्वरप्रीत्यर्थं काशी-गया-प्रयागादि-तीर्थयात्राप्रस्थानमहं करिष्ये।

इति सङ्कल्प्य, (पुनश्चाक्षतान् गृहीत्वा)

तीर्थक्षेत्रयात्राकाले सम्भाव्यमान-चोरभय-विद्युद्भय-धूर्त-मूर्ख-पाषण्ड-दस्यु-वञ्चकभय-व्याघ्रोपद्रव-भूतोपद्रव-प्रेत-पिशाच-वेताल-ब्रह्मराक्षसाद्युपद्रवनिरासार्थं
वाहनचलनकाले विपरीतगमनोन्मार्गगमन-कुमार्गप्रपतन-विमार्गगमनोत्पथप्रवृत्त्यादिदोषनिवारणार्थं
वाहनप्रकम्पनोत्कम्प-क्षेप-विक्षेप-दोषप्रतिबन्धार्थं सुखप्रयाणसिद्ध्यर्थं प्रयाणकाले
समस्तगो-विप्र-प्रसादसिद्ध्यर्थं सकलदेवतानुग्रहप्राप्त्यर्थं सकलब्राह्मणजन-ऋषिजन-बहुश्रुतब्राह्मणजनानां
शुभाशीर्वादप्राप्त्यर्थं कुलदेवताप्रीत्यर्थम् इष्टदेवताप्रसादसिद्ध्यर्थं च सकलदेवानां
ग्रहाणां च आनुकूल्यसिद्ध्यर्थं च (विप्रजनपूजनं) नान्दीमुखपितृपरितोषणं (च) करिष्ये।
तदादौ निर्विघ्नतासिद्ध्यर्थं गणपतिपूजां करिष्ये। इति सङ्कल्प्य, तत्कृत्वा विप्रजनपूजां कुर्यात्।

{\small कूर्मपुराणे —
\fourlineindentedshloka*
{गङ्गादितीर्थेषु वसन्ति मत्स्या}
{देवालये पक्षिगणाश्च सन्ति}
{भावं ततो हृत्कमले निधाय}
{तीर्थानि सेवेत समाहितात्मा}
(पुण्यतीर्थेषु भावनयैव फलमिति तात्पर्यम्।)\par}

\sect{विप्रजन-पूजा}

चतुरो विप्रान् उदङ्मुखानुपवेश्य स्वयं प्राङ्मुख उपविश्य

\begin{center}
\fourlineindentedshloka*
{उच्चैर्गतिर्जगति सिध्यति धर्मतश्चेत्}
{तस्य प्रमा च वचनैः कृतकेतरैश्चेत्}
{तेषां प्रकाशनदशा च महीसुरैश्चेत्}
{तानन्तरेण निपतेत् क्व नु मत्प्रणामः}
\end{center}

इति नमस्कृत्य

\begin{center}
ॐ विप्रपुङ्गवेभ्यो नमः, पाद्यं समर्पयामि।\\
ॐ द्विजोत्तमेभ्यो नमः, अर्घ्यं समर्पयामि।\\
ॐ महीसुरेभ्यो नमः, आचमनीयं समर्पयामि।\\
ॐ यज्ञसूत्रशोभितेभ्यो नमः, उपवीतं समर्पयामि।\\
ॐ गन्धानुलिप्तेभ्यो नमः, गन्धं समर्पयामि।\\
ॐ षट्कर्मनिरतेभ्यो नमः, अक्षतान् समर्पयामि।\\
ॐ बहुश्रुतेभ्यो नमः, अलङ्कारान् समर्पयामि।\\
ॐ श्रुतिभूषणेभ्यो नमः, पुष्पाणि समर्पयामि।\\
ॐ शास्त्रनेत्रेभ्यो नमः, धूपम् आघ्रापयामि।\\
ॐ वेदचक्षुर्भ्यो नमः, दीपं दर्शयामि।\\
ॐ अनर्घ्यवैभवेभ्यो नमः, पुनरर्घ्यं समर्पयामि।
\end{center}

\dnsub{पादपूजा}

आधिव्याधिहरं\ldots\ अच्युतं पदम्॥ श्रीपादपद्माभ्यां नमः।
गोत्रप्रवरमुक्त्वा अभिवाद्य इत्यादि\ldots॥

ॐ अध्यात्मवित्तमेभ्यो नमः। इदं वोऽष्टार्घ्यं सम्पूर्णमस्त्विति भवन्तो ब्रुवन्तु\ldots।
ॐ हिरण्यगर्भगर्भस्थं\ldots\ प्रयच्छ मे।
यात्राकाले समस्तविप्रप्रसादसिद्ध्यर्थं नानानामगोत्रेभ्यो ब्राह्मणेभ्यो यथाशक्ति हिरण्यं
(वस्त्रं, गोदानं, सुवर्णदानं वा\ldots) सम्प्रददे श्रियै न मम। इति दत्वा, अक्षतान् दत्वा,
ॐ नवो नवो भवति\ldots\ दीर्घमायुः। इत्याशीः सम्प्रार्थ्य, गणपतये नमः प्रसन्नपूजां समर्पयामि,
प्रसन्नार्घ्यं समर्पयामि। इति प्रसन्नपूजां कृत्वा यस्य स्मृत्या\ldots\ अच्युतम्, अनेन
काश्यादितीर्थक्षेत्रयात्राप्रस्थानाङ्गत्वेन महागणाधीश्वरपूजापूर्वक-विप्रजनपूजनकर्मणा
श्रीपरमेश्वरः प्रीयताम्। ॐ तत्सत्। इति ब्रह्मार्पणं कृत्वा। प्रसादं स्वीकृत्य।

\dnsub{पुण्याहवाचनम्}

\twolineshloka*
{सम्पूज्य गन्धमाल्याद्यैर्ब्राह्मणान् स्वस्ति वाचयेत्}
{धर्मकर्मणि माङ्गल्ये सङ्ग्रामेऽद्भुतदर्शने}

पुण्याहवाचनं कृत्वा (तत्र `नान्दीमुखाः पितरः प्रीयन्ताम्' इति विशेषः) नान्दीश्राद्धं च कुर्यात्।

\centerline{॥इति काश्यादितीर्थक्षेत्रयात्राप्रस्थानविधिः॥}
